
%If \(\mu\) is a real measure on \(\R\) satisfying
%\begin{equation}
%  \int_{\R}^{} \frac{d\mu(t)}{1+t^{2}} < \infty
%,\end{equation}
%then its Stieltjes transform is the function
%\begin{equation}
%  S(z) = \int_{\R}^{} \frac{d\mu(t)}{t-z}
%.\end{equation}
%It turns out one can recover the measure by the Stieltjes inversion
%formula.

One shows that, for a self-adjoint operator \(A\colon\Cal D_A\to\Cal
H\), the function \(N_\psi\colon\C^{+}\to\C\)
\begin{equation}
  N_\psi(z) = \langle \psi,R_A(z)\psi\rangle
\end{equation}
is a Herglotz-Nevanlinna function (i.e. sends
\(N_\psi(\C^{+})\subseteq \C^{+}\) ) which additionally satisfies
\begin{equation}
  \sup_{y>0} |yN_\psi(iy)|<\infty
\end{equation}
which implies that can be written in the form
\begin{equation}
  N_\psi(z)
  = \int_{\R}^{} 
    \frac{d\omega(t)}{t- z},
\end{equation}
where the measure \(\mu\) is given by
\begin{equation}
  \mu(t_1,t_2]
  =
  \lim_{\delta\to 0^{+}}
  \lim_{\epsilon\to 0^{+}}
  \frac{1}{\pi}
  \int_{t_1+\delta}^{t_2+\delta}\Im N_\psi(t+i\epsilon)\d t
.\end{equation}
Then one constructs the unique projection-valued measure \(P\)
satisfying
\begin{equation}
  \langle \psi, P(\Omega)\psi\rangle = \mu(\Omega)
.\end{equation}
We claim that \(A=\int_{M}^{}\lambda\d P(\lambda)\).


